\section{独立八尺度障碍与全体三阶亏损}\label{ic:section}
共同尺度的解析解不能直接作为任意尺度的上界。本节不作尺度对称化，而从每个反点对的
交并误差出发，得到独立尺度的八点障碍，再把它传递到真实光晕的面积。
最终闭合的是两点包络与全体常数之间的亏损阶，三阶系数本身仍未确定。

\begin{lemma}[两点交并的误差预算]\label{ic:pair}
设 $|E|=1$，两个严格见证正方形为 $Q_1,Q_2$，边长为 $\ell_1,\ell_2$，
中心坐标差的绝对值为 $u,v$，置 $z=(\ell_1+\ell_2)/2$、$I=Q_1\cap Q_2$、$V=Q_1\cup Q_2$。
在 $u,v<z$ 时，有
\[
 |I\setminus E|+|E\setminus V|
 <1+|I|-2\alpha z
 \le G(z):=1+(z-u)(z-v)-2\alpha z.
\]
若 $u+v\ge W:=\alpha^{-1}-\alpha$，则
\[
 0<G(z)\le G(\alpha)=1-\alpha^2-\alpha(u+v)+uv.
\]
\end{lemma}
\begin{proof}
精确恒等式
$|E\cap Q_1|+|E\cap Q_2|=1+|I|-|I\setminus E|-|E\setminus V|$
与严格阈值给出第一式。每个坐标交长至多为平均边长减去中心差，故得到 $G$。
合法尺度给 $\alpha<z<\alpha^{-1}$，而 $G$ 为凸二次式。
又 $G(\alpha)-G(\alpha^{-1})=W(u+v-W)\ge0$，最大端点为 $\alpha$。
这里并未要求两个见证边长相同。
\end{proof}

\begin{theorem}[任意独立八尺度的二阶径向障碍]\label{ic:radial}
固定 $0<b<1/2$、$a=1-b$，仍记 $X_r=\{(\pm ar,\pm br),(\pm br,\pm ar)\}$，
$r_0$ 如定理~\ref{cs:sharp}。存在 $c_b,t_b>0$，使 $0<t<t_b$ 时，
任何面积一可测输入的严格光晕都不能包含半径
$r\ge r_0(1-c_bt)$ 的全部八点，即使八个见证尺度完全独立。
若 $R_8(t,b)$ 是这些八点在全体输入上的可行半径上确界，则
\[
 r_0-R_8(t,b)\asymp_b t^2.
\]
\end{theorem}
\begin{proof}
半径 $r\ge r_0$ 已由两点差体排除。若所述常数不存在，可取
$t_j\downarrow0$、$0<\sigma_j\downarrow0$，以及严格可行的八点
$r_j=r_0(1-\sigma_jt_j)$。以下略去下标并沿子列取极限。
对每个反点对，$u=2ar$、$v=2br$。由两点边界恒等式及展开，
\[
 G(\alpha)=O_b(\sigma t^2)=o(t^2),\qquad
 u+v-W=4ab\,t^2+o(t^2)>0.
\]
引理~\ref{ic:pair}因而给每个反点对
\begin{equation}\label{ic:concentration}
 |I\setminus E|+|E\setminus V|=o(t^2).
\end{equation}
写各尺度 $\ell_i=1+\lambda_i t$，合法性给
$-1<\lambda_i<1+O(t)$。对反点对的平均系数 $\mu$，
\[
 G(1+\mu t)=(\mu^2-1+2\sigma)t^2+O_b(t^3)>0.
\]
因此 $\mu\to+1$ 或 $-1$，且该反点对的两个 $\lambda_i$ 必须趋于同一端点。
取子列使四个反点对的正负相位固定。

把方向分为 $A=(a,\pm b)$ 与 $B=(b,\pm a)$。
负相位的两个方体之并，其外包矩形边长为
$1+(a-b)t+o(t),1-(a-b)t+o(t)$，正负号在另一方向交换。
若负相位同时出现于 $A,B$，两外包矩形交集面积至多
$1-2(a-b)t+o(t)$，却由~\eqref{ic:concentration}几乎包含面积一的 $E$，矛盾。
正相位的反点交集则为边长 $1\pm(a-b)t+o(t)$ 的矩形；
若两方向均有正相位，两个交集的并面积为
$1+2(a-b)t+o(t)$，却须几乎包含于 $E$，同样矛盾。
故只剩 $A$ 全负、$B$ 全正，或交换的相位。

以前者为例，取任一负 $A$ 对的并 $V_A$ 与任一正 $B$ 对的交 $I_B$。
$I_B\setminus V_A$ 至少包含两个角矩形，每个边长为
$(2a+o(1))t$ 与 $(2b+o(1))t$。于是
\[
 |I_B\setminus V_A|\ge(8ab+o(1))t^2.
\]
另一方面，它不超过 $|I_B\setminus E|+|E\setminus V_A|=o(t^2)$，矛盾。
另一相位由坐标交换处理。这排除了反例序列，故存在固定的 $c_b,t_b$。
共同尺度构造给
$r_0-R_8\le r_0-R_b=\tfrac12(1-4b^2)t^2+O_b(t^3)$，
结合刚证的下界即得匹配阶。这里没有分类固定正 $\sigma$ 的全部尺度模式。
\end{proof}

\begin{lemma}[边界八盒的鲁棒障碍]\label{ic:boxes}
固定 $b\in(0,1/2)$，令 $p_i\in X_{r_0}\subset\partial T$，
$T=\overline K_\alpha/2$。在每点取单位外法向 $n_i$ 和切向 $v_i$。
存在 $\eta,t_0>0$，使 $0<t<t_0$ 时，没有同一质量一输入的严格光晕能同时命中八盒
\[
 B_i(t)=\{p_i+u n_i+v v_i:|u|\le\eta t^2,\ |v|\le\eta t\}.
\]
\end{lemma}
\begin{proof}
若否，可取 $\eta_j\to0$、$t_j\to0$ 及每盒一个严格光晕点。
在正象限，边界方程为
$(\alpha-2x)(\alpha-2y)=2\alpha^2-1$，其切向量两坐标和为 $O_b(t)$。
因此反点差的坐标符号稳定，并有
\[
 u+v=2r_0+O_b(\eta_jt_j^2),\qquad
 uv=4ab r_0^2+O_b(\eta_jt_j^2).
\]
相应 $G(\alpha)=o(t_j^2)$，且 $u+v-W=4ab t_j^2+o(t_j^2)>0$。
此前的交并预算、端点相位和角矩形矛盾均保持不变。
若不存在某组固定 $\eta,t_0$，便可按 $\eta_j,t_j<1/j$ 选出上述序列，
故序列矛盾正好给出统一的八盒障碍。
\end{proof}

\begin{lemma}[小缺面积下侧弧内盒的传递]\label{ic:transfer}
固定 $b=9/20$，将引理~\ref{ic:boxes}的 $\eta$ 缩小到 $0<\eta\le10^{-6}$。
在对应边界点定义
\[
 P_i(t)=\{p_i-u n_i+v v_i:
 3\eta t^2/8\le u\le5\eta t^2/8,\ |v|\le\eta t/8\}.
\]
它们在小 $t$ 时包含于 $T$，每个面积为 $\eta^2t^3/16$。
若 $A_j$ 是满维紧凸体、$S_j=(A_j-A_j)/2\subset T_j$，且
\[
 |T_j|-|S_j|=o(t_j^3),\qquad |S_j|-|A_j|=o(t_j^3),
\]
则可平移 $A_j$，使所有 $P_i(t_j)$ 最终均包含于 $A_j$。
\end{lemma}
\begin{proof}
记 $g=|T|-|S|$、$\Delta=|S|-|A|$。
由引理~\ref{pb:four-normal}作平移，$h_A=h_S+o$，且奇部在四个对角法向归零。
引理~\ref{pb:local-caps}及局部奇部估计给
\begin{equation}\label{ic:local-support}
 h_A(\theta)-h_T(\theta)=o(t^2)
 \quad\text{一致于 }|\theta-\theta_k|\le t/8.
\end{equation}
这里 $g=o(t^3)$ 使 $g^{2/3}=o(t^2)$，且
$\sqrt{\Delta t}=o(t^2)$。
又由引理~\ref{pb:cap-distance}，
$h_A\ge h_T-\sqrt g-\sqrt\Delta\ge t/2-o(t)$，故 $0\in\operatorname{int}A$。

对任一 $p=p_i$，其法向角为
$\theta_k\pm t/10+O(t^2)$，严格位于上述窗口内部。
以 $n=n_i,v=v_i$ 为基，使用支撑函数的一次齐次延拓，置
\[
 f_{A,t}(s)=\frac{h_A(n+st v)-p\cdot(n+st v)}{t^2}.
\]
这是 $s$ 的凸函数。由~\eqref{ic:local-support}及
$\rho_T\to\rho_*:=1/\sqrt2$，在 $|s|\le1/100$ 上一致有
\[
 f_{A,t}(s)=f_{T,t}(s)+o(1)\longrightarrow\rho_*s^2/2.
\]
凸函数一致收敛到光滑凸函数，因而紧内区间上所有次梯度也收敛到极限导数；
这可直接由左右差商夹逼证明。
在固定 $s=\pm s_0$ 处选任意支撑点 $q_\pm\in A$。
支撑恒等式和次梯度关系于是给
\[
 \frac{(q_\pm-p)\cdot v}{t}\to\pm\rho_*s_0,
 \qquad
 \frac{(q_\pm-p)\cdot n}{t^2}\to-\rho_*s_0^2/2.
\]
取 $s_0=\sqrt\eta/16$。两点的弦切向半宽为
$(\rho_*s_0+o(1))t$，法向深度小于 $\eta t^2/4$。
而 $P_i$ 的切向半宽只有 $\eta t/8$，深度在
$[3\eta/8,5\eta/8]t^2$。
第三个顶点 $0$ 的法向深度为 $(1/\sqrt2+o(1))t$、切向坐标为 $O(t)$；
在所需 $t^2$ 深度处，与它作凸组合的权重仅为 $O_\eta(t)$，切向位移为 $O_\eta(t^2)$。
因此三角形 $\operatorname{conv}\{0,q_-,q_+\}$ 最终包含整个 $P_i$，从而 $P_i\subset A$。
有限个点可以同时处理。

最后验证内盒本身在 $T$ 内。在正象限令
$F(x,y)=1-\alpha^2-2\alpha(x+y)+4xy$。
$F(p)=0$，$\nabla F(p)=-2Nn$，其中 $N\to\sqrt2$。
对 $z=-un+\xi v_i$，
$F(p+z)=2Nu+4z_1z_2\ge2Nu-2(u^2+\xi^2)>0$
于上述盒内，且坐标仍在正象限边界分支内；其他盒由反射处理。
面积公式由两边长 $\eta t^2/4$、$\eta t/4$ 得到。
\end{proof}

\begin{theorem}[两点包络的三阶亏损及其匹配阶]\label{ic:cubic-gap}
存在绝对常数 $\kappa,t_0>0$，使二维半阶中心正方形模型满足
\[
 C_{1,\mathcal Q_c}(1-t)\le\mathcal U(1-t)-\kappa t^3
 \qquad(0<t<t_0).
\]
结合矩形下界，得到
\[
 \mathcal U(1-t)-C_{1,\mathcal Q_c}(1-t)\asymp t^3.
\]
特别地，可将 $\kappa$ 缩小到 $0<\kappa<1$，使
\[
 2\le\liminf_{t\downarrow0}\frac{C_{1,\mathcal Q_c}(1-t)-2t^2}{t^3}
 \le\limsup_{t\downarrow0}\frac{C_{1,\mathcal Q_c}(1-t)-2t^2}{t^3}
 \le\frac{10}{3}-\kappa<\frac{10}{3}.
\]
常数与起效阈值未数值化；仍未证明中间极限存在。
\end{theorem}
\begin{proof}
若不存在统一的正 $\kappa,t_0$，按 $\kappa,t_0<1/j$ 对角选取
$t_j\downarrow0$。由上确界定义，可再选质量一输入，使其严格光晕 $H_j$ 满足
$\delta_j=\mathcal U(1-t_j)-|H_j|=o(t_j^3)$；无需假设上确界达到。
置 $A_j=\overline{\operatorname{conv}H_j}$、$S_j=(A_j-A_j)/2\subset T_j$。
缺面积精确分解为
\[
 \delta_j=(|T_j|-|S_j|)+(|S_j|-|A_j|)+(|A_j|-|H_j|),
\]
三项均非负，故均为 $o(t_j^3)$。
引理~\ref{ic:transfer}给平移后 $P_i(t_j)\subset A_j$。
因 $|A_j\setminus H_j|=o(t_j^3)$，而每个盒面积为固定正系数乘 $t_j^3$，
$H_j$ 必须命中全部八盒。这与引理~\ref{ic:boxes}矛盾。
平移同时作用于输入和光晕，不改变质量或算子。
序列矛盾证明了统一的正 $\kappa,t_0$。
另一方面 $C\ge R$ 且 $\mathcal U-R=(4/3)t^3+O(t^4)$，给出匹配的上侧阶及所列夹逼。
\end{proof}
